\section{Why ApportionRegions Is the Only Defensible Choice}
\label{sec:why-apportionregions}

\subsection{The Constitutional-Extension Argument}

The B-series algorithmic redistricting program \citep{b0paper} has developed
fifteen distinct algorithm variants for drawing congressional districts.
Each is mathematically defensible.
Each is politically neutral in the sense that it contains no explicit
partisan or racial inputs.
Seven of the fifteen have been empirically validated across all fifty states.

Yet only one of them---\ar{} \citep{b11paper}---is constitutionally derived.
The others are practitioners' choices.
This distinction is the entire weight-bearing structure of the first-principles
argument: a practitioner's choice, however mathematically elegant, is a
political choice once it is mandated by a federal statute.

The reason is the compactness-proportionality paradox demonstrated in
Section~\ref{sec:introduction}.
Choosing which algorithm to use determines, in many states, which party wins
more seats.
A federal statute that mandates ``use algorithm X'' is therefore making a
political choice.
The only way to escape this logic is to mandate an algorithm whose identity
is determined by first principles that precede and transcend any practitioner
preference.
\hh{} is that algorithm.
\ar{} is its graph-partitioning realization.

\subsection{The Prime Factorization as Deterministic Bisection Tree}

The key insight in \ar{} \citep{b11paper} is that the Huntington-Hill priority
rule, applied to a state's $k$ congressional seats, determines a unique
bisection tree.
This tree tells the algorithm exactly how to recursively divide the state's
population.

The mechanism is the prime factorization of $k$.
For example:
\begin{itemize}
  \item California ($k = 52 = 4 \times 13$): The HH priority rule dictates
    a $4$-way first-level split (into four regions of 13 districts each),
    followed by thirteen $13$-way splits.
  \item Texas ($k = 38$): $38 = 2 \times 19$.
    First split: $2$-way (two regions of 19).
    Second: $19$-way (each region into 19 districts).
  \item New York ($k = 26 = 2 \times 13$): First split $2$-way, then $13$-way.
  \item Oregon ($k = 6 = 2 \times 3$): First split $2$-way (East/West),
    then each half $3$-way.
\end{itemize}

The prime factorization is unique (fundamental theorem of arithmetic).
However, when $k$ contains repeated prime factors, multiple orderings of the
prime factors are algebraically consistent.
For example, $k = 12 = 2^2 \times 3$ admits two factorisation orderings:
$[3, 2, 2]$ (three-way first split, then two-way twice) and $[2, 3, 2]$ and $[2, 2, 3]$
(two-way first split in different arrangements).
The \dia{} statute resolves this ambiguity by mandating a \emph{canonical ordering}:
prime factors are applied in \textbf{non-increasing order} (largest prime factor first,
then descending).
Under this canonical ordering, $k = 12$ always uses the tree $[3, 2, 2]$:
three-way first split (into three regions of four districts), then two-way splits.
For New Jersey ($k = 12$), this canonical ordering is the only valid \ar{} tree.\footnote{%
  The canonical ordering---largest prime factor first---also maximises
  early-stage geographic compactness.
  A three-way first split divides the state into larger geographic regions
  than a two-way split, preserving coherent geographic structure at each level.
  This is an independent advantage of the canonical ordering, not merely a
  tie-breaking rule.
}

Given $k$ and the canonical ordering, there is exactly one \hh{} bisection tree.
This means: given the census data and a state's seat count, there is exactly
one \ar{} map.
Not one of many acceptable maps.
One map.

This uniqueness is the source of the algorithm's constitutional legitimacy.
A court or outside reviewer could verify that a named implementation produced
a named output from published inputs and parameters. That does not make the map
mathematically unique: structure, geographic units, adjacency rules, engine,
seed, weights, and exception procedures must first be fixed by law or rule.

\subsection{Seed Sensitivity: A Required Empirical Test}

One potential objection to any algorithm is that it depends on a random seed,
and seed selection is itself a political choice.
\citet{b11paper} motivates testing whether \ar{} has low seed sensitivity.
The available NC and WI experiments are encouraging, but they do not establish
that the resulting map is identical across seeds nationally.

This result is a consequence of the tight geographic structure of the
\hh{} bisection tree.
The prime-factor tree can constrain the candidate space, but METIS remains a
heuristic local optimizer. Companion paper T.4 reports limited state examples;
a committed 50-state seed-trace package is still required before making a
national variance claim.

A content-derived seed can prevent post hoc seed shopping when its schema and
all seed-affecting inputs are precommitted. It does not make census data the
only influence on the map and does not eliminate practitioner or legislative
choices. The current canonical candidate derives the seed from the complete
assignment-affecting manifest rather than the census release identifier alone.

\subsection{Seed-Sweep Evidence And Remaining Package Gap}

\citet{b7paper} reports a 50-seed comparison for the edge-weighted bisection
family, including seat-count standard deviations below 0.3 in the reported
states. That paper-level result is not a substitute for a committed 50-state
seed-trace package, and it does not certify \ar{} map identity or convergence
across the full solution space.
The minimum-edge-cut objective is a strong geographic signal, but the relative
effect of initialization must be measured rather than assumed.

For \ar{}, prime factorization fixes a tree-ordering rule once its canonical
ordering and prime-seat fallback are selected. It does not prove that METIS
returns the same geographic partition or that seed variation is confined to
individual boundary tracts.

\subsection{The Other Algorithms Are Practitioners' Choices}

To understand why \ar{} is uniquely defensible, it helps to understand
why the other algorithms in the toolbox are not.

\paragraph{GeoSection (T.1).}
GeoSection \citep{b8paper} selects the bisection ratio by minimizing the
isoperimetrically normalized edge-cut $\mathrm{EC}_\text{norm}(i:j) = \mathrm{EC}(i:j)/\sqrt{\min(i,j)}$.
This is a mathematically elegant criterion---it selects the split that
produces the most compact division relative to the expected compactness
at each scale.
But a practitioner must choose to use GeoSection rather than some other
ratio-selection criterion.
That choice, as the Wisconsin bakeoff shows, determines whether Democrats
win four seats or three.
GeoSection is a rational choice.
It is not a constitutionally derived choice.

\paragraph{AreaSection (T.2).}
AreaSection \citep{b9paper} balances population and land area simultaneously,
giving rural areas somewhat larger districts by geography and urban areas
somewhat smaller districts by geography.
This is a coherent policy choice: one might argue that rural areas have
different representation needs.
But it is a policy choice, not a constitutional derivation.
The Constitution does not say anything about land area.
\wesberry{} is explicit: it is population, not geography, that determines
the value of a vote.
AreaSection would systematically advantage rural (Republican) areas.
A statute mandating AreaSection would be a partisan choice wearing an
algorithmic costume.

\paragraph{Unweighted bisection (B.1).}
The simplest algorithm---bisect the state's population into two equal halves
at each level, ignoring geographic boundary lengths---produces compact
districts in some states and bizarre districts in others.
In Wisconsin, it produces a 4D/4R outcome.
In California, it splits the Central Valley in ways that create non-compact
districts.
The absence of geographic weighting is itself a practitioner's choice.
One could argue for it---simplicity has constitutional value---but one must
argue for it.
The Constitution does not mandate simplicity.

\paragraph{Why none of the others is constitutionally derived.}
All of GeoSection, AreaSection, unweighted bisection, and the eleven other
B-series variants require a practitioner to choose them.
Each choice is defensible.
None is compelled.
When Congress mandates one of them, it is choosing to produce certain
electoral outcomes, because algorithm selection determines electoral outcomes,
as the Wisconsin bakeoff proves.

\ar{} is the sole exception because its selection criterion---\hh{} applied
to the intrastate problem---is derived from the same principle that Congress
already mandated for the interstate problem.
Choosing \ar{} is not choosing a political outcome.
It is applying a previously enacted constitutional standard at the next level
of the hierarchy.

\subsection{The Compactness-Proportionality Paradox as Legal Doctrine}

We formalize the Wisconsin observation into a doctrinal test.

\begin{definition}[Practitioner's Algorithm]
An algorithm is a \emph{practitioner's algorithm} if there exists a
pair of states $(A, B)$ and a pair of different mandated algorithms
$(\mathcal{A}_1, \mathcal{A}_2)$ such that $\mathcal{A}_1$ and $\mathcal{A}_2$
produce materially different seat-count outcomes in at least one state,
and neither outcome can be derived from the constitutional text without
reference to $\mathcal{A}_1$ or $\mathcal{A}_2$.
\end{definition}

\begin{definition}[Constitutionally Derived Algorithm]
An algorithm is \emph{constitutionally derived} if its selection criterion
is uniquely determined by an existing constitutional or statutory mandate,
without any practitioner choice at the level of algorithmic selection.
\end{definition}

\begin{proposition}
GeoSection, AreaSection, unweighted bisection, and all other B-series
algorithms except \ar{} are practitioners' algorithms.
\ar{} is constitutionally derived.
\end{proposition}

\begin{proof}
By the Wisconsin bakeoff \citep{b0paper}: GeoSection and unweighted bisection
produce materially different seat-count outcomes in Wisconsin
(3D/5R vs.\ 4D/4R, a 12.5-percentage-point partisan gap).
Neither outcome is derived from the constitutional text without reference
to the choice of algorithm.
Hence both are practitioners' algorithms.
The argument extends to all other non-\ar{} algorithms by the same
Wisconsin evidence.

For \ar{}: its selection criterion is the Huntington-Hill priority rule.
The Huntington-Hill priority rule is uniquely mandated by 2~U.S.C.~\usca{}~2a(a)
for the interstate problem.
The intrastate and interstate problems share the same decision structure
(Section~\ref{sec:first-principles}): the \hh{} priority sequence determines
the bisection tree, and the bisection tree is derived without practitioner choice
from the prime factorization of $k$.
The adjacency constraint does not introduce a practitioner choice at the level
of tree selection---it constrains the feasible set of realizations within which
METIS finds the minimum-edge-cut partition.
Therefore, applying \hh{} at the intrastate level to determine the bisection tree
requires no additional practitioner choice beyond the choice already embodied in
federal statute; the contiguous realization is determined by the METIS minimum-edge-cut
objective, which is itself free of partisan input.
\ar{} is uniquely constitutionally derived.
\qed
\end{proof}

The practical implication of this proposition is significant.
A federal statute mandating any algorithm other than \ar{} can be challenged
as an implicit political choice, because every other algorithm produces
politically different outcomes, and its selection cannot be grounded in
existing constitutional doctrine.
A federal statute mandating \ar{} cannot be so challenged, because \ar{}
is not a practitioner's choice---it is the application of an existing
standard at the next level of the hierarchy.

\subsection{The Factorization Spine and Multi-Chamber Compatibility}

A subsidiary advantage of the \hh{}-derived bisection tree is compatibility
with multi-chamber redistricting.
\citet{b13paper} analyzes the compatibility of congressional and state-legislative
factorization spines, finding that 11 states have spine score 0 (exact
nesting): Oregon, Alabama, Idaho, Maine, Montana, Nebraska, Nevada,
New Hampshire, North Dakota, Rhode Island, and South Dakota.

Nesting means that the congressional and state-legislative district boundaries
can be made exactly compatible: every state legislative district is contained
within exactly one congressional district, or every congressional district
contains an exact integer number of state legislative districts.
This is a governance benefit: constituents know exactly which legislative
bodies have jurisdiction over which geographic areas, and no census tract
straddles the boundary between a state-legislative and congressional district.

This benefit is not available with practitioners' algorithms.
When a practitioner chooses GeoSection for congressional maps and
a different algorithm for state legislative maps, the resulting boundaries
may bear no rational relationship to each other.
The \hh{} factorization spine produces coherent multi-chamber hierarchies
as a byproduct of first-principles derivation.
\noindent\textbf{SUPERSEDED SOURCE --- NOT COMPILED.} This section belongs to
the former ApportionRegions final-map mandate. See
\texttt{07-canonical-alignment.tex} and \texttt{REVISION-PLAN.md}.
